\documentclass[11pt]{article}

\usepackage[margin=1in]{geometry}
\usepackage{amsmath,amssymb,mathtools}
\usepackage{bm}
\usepackage{booktabs}
\usepackage{array}
\usepackage{enumitem}
\usepackage{microtype}
\usepackage[colorlinks=true,linkcolor=blue,urlcolor=blue,citecolor=blue]{hyperref}

\newcommand{\ii}{\mathrm{i}}
\newcommand{\ee}{\mathrm{e}}
\newcommand{\cQ}{\mathcal{Q}}
\newcommand{\cC}{\mathcal{C}}
\newcommand{\cL}{\mathcal{L}}
\newcommand{\cD}{\mathcal{D}}
\newcommand{\ip}[2]{\left\langle #1,#2\right\rangle}
\newcommand{\cc}{\text{c.c.}}
\newcommand{\T}{\mathbb{T}}
\newcolumntype{L}[1]{>{\raggedright\arraybackslash}p{#1}}
\DeclareMathOperator{\RePart}{Re}
\DeclareMathOperator{\spanop}{span}

\title{Multiple-Scales Computations for Proposal 3}
\author{}
\date{\today}

\begin{document}

\maketitle
\tableofcontents

\section{Purpose and scope}

This note carries out the general calculation proposed for a one-dimensional,
translation- and reflection-invariant pattern-forming system near a finite
wavenumber Hopf (Turing--Hopf) bifurcation.  The calculation is fully explicit
in terms of the Fourier matrices of the linear operator and the quadratic and
cubic multilinear maps.  Numerical values cannot be produced until the
specific matrices, kernels, and forcing vectors are supplied.

Two corrections to the handwritten derivation are important from the outset.

\begin{enumerate}[label=\arabic*.]
  \item The modes with spatial wavenumbers $+k_0$ and $-k_0$ have independent
  complex amplitudes.  They must be denoted, for example, by $A(\tau)$ and
  $B(\tau)$.  Setting them equal in advance restricts the calculation to a
  standing-wave subspace and makes it impossible to decide between standing
  and traveling waves.

  \item If the slow time is $\tau=\epsilon^2t$, the chain rule is
  \begin{equation}
    \partial_t=\partial_s+\epsilon^2\partial_\tau,
    \qquad s=t,
    \label{eq:chain-rule}
  \end{equation}
  not $\partial_s+\epsilon\partial_\tau$.  The latter belongs to the different
  scaling $\tau=\epsilon t$.
\end{enumerate}

\section{Equation, notation, and hypotheses}

\subsection{Corrected notation}

The proposal uses $Q$ both for the quadratic nonlinearity and for a forcing
vector.  To avoid that collision, write the equation as
\begin{align}
  u_t={}&\cL_0u+p\cL_1u+\cQ(u,u)+\cC(u,u,u) \notag\\
       &+\epsilon^n\left(P\ee^{\ii\omega t}+\overline P\ee^{-\ii\omega t}\right)
       +\epsilon^m\left(S\ee^{\ii kx}+\overline S\ee^{-\ii kx}\right),
  \label{eq:pde}
\end{align}
where
\begin{itemize}
  \item $u(x,t)\in\mathbb{R}^d$ and $x\in\T:=\mathbb{R}/(2\pi\mathbb{Z})$;
  \item $p$ is the bifurcation parameter (called $\lambda$ in the proposal);
  \item $\cQ$ is a symmetric bilinear map and $\cC$ is a symmetric trilinear
  map, both extended complex linearly during the calculation;
  \item $P,S\in\mathbb{C}^d$ are the temporal- and spatial-forcing vectors.
\end{itemize}
Any quadratic or cubic form has a symmetric multilinear polarization, so the
symmetry assumptions on $\cQ$ and $\cC$ lose no generality.  If an implementation
stores a nonsymmetric bilinear map, every occurrence of $2\cQ(a,b)$ below must
be replaced by $\cQ(a,b)+\cQ(b,a)$.

For each integer spatial Fourier mode $h$, define matrices $M_j(h)$ by
\begin{equation}
  \cL_j\!\left(v\ee^{\ii hx}\right)
  =M_j(h)v\ee^{\ii hx},
  \qquad j=0,1.
  \label{eq:fourier-symbol}
\end{equation}
\paragraph{Meaning and construction of $M_j(h)$.}
On the $2\pi$-periodic domain, a vector field has the Fourier expansion
\begin{equation}
  u(x,t)=\sum_{h\in\mathbb Z}\widehat u_h(t)\ee^{\ii hx},
  \qquad \widehat u_h(t)\in\mathbb C^d.
\end{equation}
Translation invariance prevents the linear operator from mixing different
values of $h$.  It can change only the $d$ component amplitudes of a pure
mode, so its action on that mode is represented by the $d\times d$ matrix
$M_j(h)$.  Consequently
\begin{equation}
  \cL_j u=\sum_{h\in\mathbb Z}M_j(h)\widehat u_h\ee^{\ii hx}.
\end{equation}
The matrices are obtained by applying $\cL_j$ to $v\ee^{\ii hx}$ and using
the Fourier replacement rules
\begin{equation}
  \partial_x\mapsto \ii h,\qquad
  \partial_{xx}\mapsto-h^2,\qquad
  W*\mapsto\widehat W(h).
\end{equation}
For example, with
$\widehat W(h)=\int_0^{2\pi}W(z)\ee^{-\ii hz}\,dz$,
\begin{equation}
  W*\left(a\ee^{\ii h\,\cdot}\right)(x)
  =\widehat W(h)a\ee^{\ii hx}.
\end{equation}
Thus a differential operator $A+D\partial_{xx}+C\partial_x$ has symbol
$A-h^2D+\ii hC$, while a neural-field convolution contributes the appropriate
$\widehat W_{ab}(h)$ to each population-to-population matrix entry.  If
$\cL(p)=\cL_0+p\cL_1+O(p^2)$, then similarly
$M(h,p)=M_0(h)+pM_1(h)+O(p^2)$, so $M_1(h)=\partial_pM(h,0)$.

Because the original operator is real, complex conjugating
\eqref{eq:fourier-symbol} gives
$M_j(-h)=\overline{M_j(h)}$.  If it also commutes with reflection
$x\mapsto-x$, then $M_j(-h)=M_j(h)$; for neural fields this usually follows
from even coupling kernels.  Hence the $+h$ and $-h$ modes have the same
linear spectrum, although their amplitudes remain independent until a
standing-wave restriction is imposed.

\subsection{Finite-wavenumber Hopf hypothesis}

Let $k_0\in\mathbb{Z}\setminus\{0\}$ and $\omega_0>0$.  Assume
\begin{equation}
  M_0(k_0)\Phi=\ii\omega_0\Phi,
  \qquad
  M_0(k_0)^*\Psi=-\ii\omega_0\Psi,
  \qquad
  \Psi^*\Phi=1.
  \label{eq:eigenvectors}
\end{equation}
Thus $\Phi$ is a right eigenvector and $\Psi$ is the adjoint/left
eigenvector.  The star is the conjugate transpose; a plain transpose is not
sufficient after complexifying the problem.

Let $\nu(h,p)$ denote the eigenvalue branch through
$\nu(k_0,0)=\ii\omega_0$.  Standard eigenvalue perturbation gives
\begin{equation}
  \partial_p\nu(k_0,0)=\Psi^*M_1(k_0)\Phi.
  \label{eq:eigenvalue-derivative}
\end{equation}
The transversality condition remembered as ``$d\lambda/dp>0$'' is more
precisely
\begin{equation}
  \RePart\!\left(\Psi^*M_1(k_0)\Phi\right)>0.
  \label{eq:transversality}
\end{equation}
The imaginary part only changes the oscillation frequency to leading order.

We also assume that there are no accidental quadratic resonances: the matrices
\begin{equation}
  2\ii\omega_0I-M_0(0),\quad
  2\ii\omega_0I-M_0(\pm2k_0),\quad
  -M_0(0),\quad -M_0(\pm2k_0)
  \label{eq:nonresonance-matrices}
\end{equation}
are invertible.  If one is singular, the center space is larger and the
two-amplitude normal form derived here is not closed.

\section{Scaling and the three order equations}

Use the weakly nonlinear Hopf scaling
\begin{equation}
  p=\epsilon^2\mu,
  \qquad s=t,
  \qquad \tau=\epsilon^2t,
  \qquad
  u=\epsilon u_1+\epsilon^2u_2+\epsilon^3u_3+O(\epsilon^4).
  \label{eq:scaling}
\end{equation}
If one instead writes $p=\epsilon p_1+\epsilon^2p_2+\cdots$, consistency with
the slow time $\tau=\epsilon^2t$ requires $p_1=0$ and $p_2=\mu$.  A nonzero
$p_1$ generally creates resonant evolution on the faster scale
$\tau=\epsilon t$.

Define
\begin{equation}
  \cD:=\partial_s-\cL_0
  \label{eq:D-definition}
\end{equation}
and let $F_j$ denote the part of the two external forcings whose exponent is
$\epsilon^j$:
\begin{align}
  F_j={}&\bm{1}_{\{n=j\}}
  \left(P\ee^{\ii\omega s}+\overline P\ee^{-\ii\omega s}\right)\notag\\
  &+\bm{1}_{\{m=j\}}
  \left(S\ee^{\ii kx}+\overline S\ee^{-\ii kx}\right).
  \label{eq:Fj}
\end{align}
Substitution of \eqref{eq:scaling} into \eqref{eq:pde}, followed by collecting
powers of $\epsilon$, gives the complete first three equations:
\begin{subequations}
\label{eq:moment-equations}
\begin{align}
  O(\epsilon):\quad
  &\cD u_1=F_1,
  \label{eq:o1}\\
  O(\epsilon^2):\quad
  &\cD u_2=\cQ(u_1,u_1)+F_2,
  \label{eq:o2}\\
  O(\epsilon^3):\quad
  &\cD u_3=-\partial_\tau u_1+\mu\cL_1u_1
  +2\cQ(u_1,u_2)+\cC(u_1,u_1,u_1)+F_3.
  \label{eq:o3}
\end{align}
\end{subequations}
These are order equations, sometimes informally called the first, second, and
third ``moment'' equations.  Sections \ref{sec:u1}--\ref{sec:normal-form}
first derive the unforced normal form.  Section \ref{sec:forcing} then restores
every forcing term.

\section{The \texorpdfstring{$O(\epsilon)$}{O(epsilon)} solution}
\label{sec:u1}

Temporarily put $F_1=0$.  Introduce the fast phases
\begin{equation}
  \xi_+=\omega_0s+k_0x,
  \qquad
  \xi_-=\omega_0s-k_0x,
  \qquad
  E_\pm=\ee^{\ii\xi_\pm}.
  \label{eq:phases}
\end{equation}
The four-dimensional real kernel of $\cD$ is
\begin{equation}
  \ker\cD=spanop_{\mathbb{R}}
  \left\{\Phi E_+,\Phi E_-,\overline\Phi\,\overline E_+,
  \overline\Phi\,\overline E_-\right\}.
\end{equation}
Consequently the most general real leading-order solution is
\begin{equation}
  u_1=A(\tau)\Phi E_+ +B(\tau)\Phi E_-+\cc
  \label{eq:u1}
\end{equation}
or, with the conjugates shown,
\begin{equation}
  u_1=A\Phi E_+ +B\Phi E_-
  +\overline A\,\overline\Phi\,\overline E_+
  +\overline B\,\overline\Phi\,\overline E_-.
  \label{eq:u1-full}
\end{equation}
The amplitudes $A$ and $B$ vary only on the slow time $\tau$.  The choice
$A=B$ is a standing-wave restriction, while $B=0$ or $A=0$ produces one of
the two traveling-wave directions.

\section{The complete \texorpdfstring{$O(\epsilon^2)$}{O(epsilon squared)} calculation}
\label{sec:u2}

\subsection{Why the \texorpdfstring{$O(\epsilon)$}{O(epsilon)} parameter
coefficient vanishes}

\begin{equation}
  p=\epsilon p_1+\epsilon^2p_2+\cdots,
  \qquad p_1,p_2\in\mathbb{R}.
  \label{eq:p-preliminary-scaling}
\end{equation}
Before imposing the distinguished scaling in \eqref{eq:scaling}, take
\eqref{eq:p-preliminary-scaling} as the more general real-parameter
expansion.  The linear operator and its spatial Fourier symbol expand as
\begin{equation}
  \cL(p)=\cL_0+p\cL_1+O(p^2),
  \qquad
  M(h,p)=M_0(h)+pM_1(h)+O(p^2).
  \label{eq:L-preliminary-expansion}
\end{equation}
Since $u=\epsilon u_1+\epsilon^2u_2+\cdots$, these expansions give
\begin{equation}
  p\cL_1u=\epsilon^2p_1\cL_1u_1+O(\epsilon^3).
\end{equation}
Thus, temporarily setting the external forcing to zero as in the unforced
normal-form calculation, the second-order equation would be
\begin{equation}
  \cD u_2=\cQ(u_1,u_1)+p_1\cL_1u_1.
  \label{eq:o2-preliminary}
\end{equation}

The four critical spatio-temporal waves in \eqref{eq:u1-full} are
\begin{equation}
  \Phi E_+,\qquad \Phi E_-,\qquad
  \overline\Phi\,\overline E_+,\qquad
  \overline\Phi\,\overline E_-.
\end{equation}
Translation invariance makes $\cL_1$ diagonal in spatial Fourier modes, so
\begin{align}
  \cL_1u_1={}&A M_1(k_0)\Phi E_+
  +B M_1(-k_0)\Phi E_- \notag\\
  &+\overline A\,\overline{M_1(k_0)\Phi}\,\overline E_+
  +\overline B\,\overline{M_1(-k_0)\Phi}\,\overline E_-.
  \label{eq:L1-u1-four-waves}
\end{align}
Every term in \eqref{eq:L1-u1-four-waves} has one of the four critical
phases.  Hence $p_1\cL_1u_1$ is resonant with $\cD$.

The corresponding adjoint critical waves are
\begin{equation}
  V_+=\Psi E_+,\qquad V_-=\Psi E_-,\qquad
  \overline V_+=\overline\Psi\,\overline E_+,\qquad
  \overline V_-=\overline\Psi\,\overline E_-.
\end{equation}
Apply the Fredholm alternative to \eqref{eq:o2-preliminary}, using the
normalized double integral \eqref{eq:inner-product}.  The expansion below
shows that $\cQ(u_1,u_1)$ contains no critical phase, while distinct
space--time Fourier harmonics are orthogonal.  Projection onto the two
independent complex adjoint modes and their conjugates therefore gives
\begin{subequations}
\label{eq:o2-four-projections}
\begin{align}
  0&=\ip{V_+}{\cQ(u_1,u_1)+p_1\cL_1u_1}
    =p_1A\alpha_+,
  &\alpha_+&:=\Psi^*M_1(k_0)\Phi,\\
  0&=\ip{V_-}{\cQ(u_1,u_1)+p_1\cL_1u_1}
    =p_1B\alpha_-,
  &\alpha_-&:=\Psi^*M_1(-k_0)\Phi,\\
  0&=\ip{\overline V_+}{\cQ(u_1,u_1)+p_1\cL_1u_1}
    =p_1\overline A\,\overline{\alpha_+},&&\\
  0&=\ip{\overline V_-}{\cQ(u_1,u_1)+p_1\cL_1u_1}
    =p_1\overline B\,\overline{\alpha_-}.&&
\end{align}
\end{subequations}
The last two conditions are the conjugates of the first two, so there are two
independent complex solvability conditions, equivalently four real ones.

These Fredholm coefficients are exactly the derivatives of the two critical
eigenvalue branches.  For $\sigma\in\{+,-\}$, put $h_\sigma=\sigma k_0$ and
let
\begin{equation}
  M(h_\sigma,p)\Phi_\sigma(p)
  =\nu_\sigma(p)\Phi_\sigma(p),
  \qquad
  \nu_\sigma(0)=\ii\omega_0,\qquad
  \Phi_\sigma(0)=\Phi.
\end{equation}
Reflection symmetry permits the same normalized adjoint eigenvector $\Psi$
at $h_+=k_0$ and $h_-=-k_0$.  Differentiating at $p=0$ yields
\begin{equation}
  M_1(h_\sigma)\Phi+M_0(h_\sigma)\Phi_\sigma'(0)
  =\nu_\sigma'(0)\Phi+\ii\omega_0\Phi_\sigma'(0).
\end{equation}
Left multiplication by $\Psi^*$ cancels the terms containing
$\Phi_\sigma'(0)$ because
$\Psi^*M_0(h_\sigma)=\ii\omega_0\Psi^*$, and $\Psi^*\Phi=1$ gives
\begin{equation}
  \boxed{
  \alpha_\pm=\Psi^*M_1(\pm k_0)\Phi
  =\partial_p\nu_\pm(0).}
  \label{eq:alpha-eigenvalue-derivatives}
\end{equation}
Here $M_1(\pm k_0)$ is the file's notation for the Fourier symbol often
written $\widehat{\cL}_1(\pm k_0)$.

By reflection symmetry, $\nu_+(p)=\nu_-(p)$ and
$\alpha_+=\alpha_-$.  The transverse-crossing assumption
\eqref{eq:transversality} therefore implies
\begin{equation}
  \RePart(\alpha_\pm)>0,
  \qquad\text{in particular}\qquad \alpha_\pm\ne0.
\end{equation}
If the leading critical solution is nontrivial, then $(A,B)\ne(0,0)$.
At least one of the first two conditions in \eqref{eq:o2-four-projections}
then forces
\begin{equation}
  \boxed{p_1=0.}
  \label{eq:p1-zero}
\end{equation}
Thus the $O(\epsilon)$ parameter displacement is not compatible with a
nonzero $O(\epsilon)$ critical amplitude on the slow scale
$\tau=\epsilon^2t$.  From this point onward one must use
\begin{equation}
  p=\epsilon^2p_2+\cdots=\epsilon^2\mu+\cdots,
  \label{eq:p-distinguished-after-solvability}
\end{equation}
so the parameter term first enters the resonant solvability condition at
$O(\epsilon^3)$.  With $p_1=0$, \eqref{eq:o2-preliminary} reduces to the
unforced version of \eqref{eq:o2}, and the second-order calculation proceeds
as follows.

Define the two basic quadratic vectors
\begin{equation}
  q_{20}:=\cQ(\Phi,\Phi),
  \qquad
  q_{11}:=\cQ(\Phi,\overline\Phi).
  \label{eq:q20q11}
\end{equation}
For a real symmetric quadratic form, $q_{11}$ is real.  Direct multiplication
of all four terms in \eqref{eq:u1-full} gives
\begin{align}
\cQ(u_1,u_1)={}&q_{20}
  \left(A^2E_+^2+B^2E_-^2+2AB E_+E_-\right)+\cc\notag\\
&+2q_{11}\left(
  |A|^2+|B|^2
  +A\overline B E_+\overline E_-
  +\overline A B\overline E_+E_-
  \right).
\label{eq:Q-u1-expanded}
\end{align}
The phase products are
\begin{equation}
\begin{gathered}
  E_+^2=\ee^{2\ii(\omega_0s+k_0x)},\qquad
  E_-^2=\ee^{2\ii(\omega_0s-k_0x)},\\
  E_+E_-=\ee^{2\ii\omega_0s},\qquad
  E_+\overline E_-=\ee^{2\ii k_0x},\qquad
  \overline E_+E_-=\ee^{-2\ii k_0x}.
\end{gathered}
\label{eq:phase-products}
\end{equation}
Thus the quadratic nonlinearity generates precisely the harmonics
\begin{equation}
  (\Omega,h)=(\pm2\omega_0,\pm2k_0),\quad
  (\pm2\omega_0,0),\quad (0,\pm2k_0),\quad (0,0).
\end{equation}
None is critical under the nonresonance hypothesis.

For later use define the resolvent
\begin{equation}
  R(\Omega,h):=\left(\ii\Omega I-M_0(h)\right)^{-1}.
  \label{eq:resolvent}
\end{equation}
Define the second-order response vectors
\begin{subequations}
\label{eq:h-vectors}
\begin{align}
  h_{20}^{+}&=R(2\omega_0,2k_0)q_{20},
  &h_{20}^{-}&=R(2\omega_0,-2k_0)q_{20},\\
  h_{20}^{0}&=2R(2\omega_0,0)q_{20},
  &h_{11}^{0}&=2R(0,0)q_{11},\\
  h_{11}^{+}&=2R(0,2k_0)q_{11},
  &h_{11}^{-}&=2R(0,-2k_0)q_{11}.
\end{align}
\end{subequations}
The factors of two in \eqref{eq:h-vectors} come directly from
\eqref{eq:Q-u1-expanded}.  Since
\begin{equation}
  \cD\left(v\ee^{\ii(\Omega s+hx)}\right)
  =\left(\ii\Omega I-M_0(h)\right)v\ee^{\ii(\Omega s+hx)},
  \label{eq:D-mode}
\end{equation}
a particular solution of \eqref{eq:o2}, with no kernel component added, is
\begin{align}
u_2={}&A^2h_{20}^{+}E_+^2+B^2h_{20}^{-}E_-^2
       +ABh_{20}^{0}E_+E_-+\cc\notag\\
&+(|A|^2+|B|^2)h_{11}^{0}
  +A\overline B h_{11}^{+}E_+\overline E_-
  +\overline A B h_{11}^{-}\overline E_+E_-.
\label{eq:u2}
\end{align}
Reality and reflection symmetry imply the needed conjugacy relations among
$h_{20}^{\pm}$ and $h_{11}^{\pm}$.  A homogeneous kernel contribution to
$u_2$ can be absorbed into a near-identity redefinition of $A$ and $B$, so the
choice \eqref{eq:u2} is the standard normalization.

\section{Fredholm alternative and the adjoint double integral}
\label{sec:fredholm}

Let
\begin{equation}
  T_0=\frac{2\pi}{\omega_0}
\end{equation}
and use the normalized space--time inner product
\begin{equation}
  \ip{v}{w}
  :=\frac{1}{2\pi T_0}
  \int_0^{T_0}\int_0^{2\pi}v(s,x)^*w(s,x)\,dx\,ds.
  \label{eq:inner-product}
\end{equation}
Periodic integration by parts shows that
\begin{equation}
  \cD^*=-\partial_s-\cL_0^*.
  \label{eq:adjoint-D}
\end{equation}
Indeed, the adjoint critical functions
\begin{equation}
  V_+=\Psi E_+,
  \qquad V_-=\Psi E_-
  \label{eq:adjoint-modes}
\end{equation}
satisfy $\cD^*V_\pm=0$ because
\begin{equation}
  \left(-\ii\omega_0I-M_0(k_0)^*\right)\Psi=0.
\end{equation}
If one defines the leading operator with the opposite sign,
$\widetilde\cD=\cL_0-\partial_s$, its adjoint also changes by an overall sign.
This explains handwritten variants resembling
$-\partial_s+M_0(k_0)^T$; all variants give the same solvability condition when
used consistently.

The Fredholm alternative says that
\begin{equation}
  \cD w=f
\end{equation}
has a $T_0$-periodic-in-$s$, $2\pi$-periodic-in-$x$ solution only if $f$ is
orthogonal to every element of $\ker\cD^*$.  At third order this means
\begin{subequations}
\label{eq:fredholm-integrals}
\begin{align}
0={}&\frac{1}{2\pi T_0}
\int_0^{T_0}\int_0^{2\pi}
\Psi^*\ee^{-\ii(\omega_0s+k_0x)}F^{(3)}(s,x)\,dx\,ds,
\label{eq:fredholm-plus}\\
0={}&\frac{1}{2\pi T_0}
\int_0^{T_0}\int_0^{2\pi}
\Psi^*\ee^{-\ii(\omega_0s-k_0x)}F^{(3)}(s,x)\,dx\,ds,
\label{eq:fredholm-minus}
\end{align}
\end{subequations}
where, without forcing for the moment,
\begin{equation}
  F^{(3)}=-\partial_\tau u_1+\mu\cL_1u_1
  +2\cQ(u_1,u_2)+\cC(u_1,u_1,u_1).
  \label{eq:F3-unforced}
\end{equation}
The two integrals simply extract the coefficients of $E_+$ and $E_-$ and
then left-multiply those vector coefficients by $\Psi^*$.

\section{The complete resonant \texorpdfstring{$O(\epsilon^3)$}{O(epsilon cubed)} terms}
\label{sec:o3-resonance}

Only terms proportional to $E_+$ enter \eqref{eq:fredholm-plus}; terms
proportional to $E_-$ enter \eqref{eq:fredholm-minus}.  Orthogonality of all
other Fourier harmonics makes them integrate to zero.

\subsection{Slow derivative and parameter terms}

The $E_+$ coefficient is
\begin{equation}
  \left[-\partial_\tau u_1\right]_{E_+}=-A_\tau\Phi,
  \qquad
  \left[\mu\cL_1u_1\right]_{E_+}
  =\mu A M_1(k_0)\Phi.
  \label{eq:linear-resonant}
\end{equation}
The corresponding $E_-$ terms are obtained by replacing $A$ by $B$.

\subsection{Cubic term}

To obtain $E_+$ from three first-order phases, the allowed products are
\begin{equation}
  E_+E_+\overline E_+=E_+,
  \qquad
  E_+E_-\overline E_-=E_+.
\end{equation}
There are three permutations of the first product and six permutations of
the second.  Hence
\begin{align}
\left[\cC(u_1,u_1,u_1)\right]_{E_+}
={}&3A|A|^2\cC(\Phi,\Phi,\overline\Phi)\notag\\
&+6A|B|^2\cC(\Phi,\Phi,\overline\Phi).
\label{eq:C-resonant}
\end{align}

\subsection{Quadratic feedback through \texorpdfstring{$u_2$}{u2}}

The complete list of phase pairings from $u_1u_2$ that produce $E_+$ is
\begin{align*}
E_+\cdot1&=E_+,
&\overline E_+\cdot E_+^2&=E_+,\\
\overline E_-\cdot(E_+E_-)&=E_+,
&E_-\cdot(E_+\overline E_-)&=E_+.
\end{align*}
Using \eqref{eq:u2}, these four pairings give
\begin{align}
\left[2\cQ(u_1,u_2)\right]_{E_+}
={}&2A(|A|^2+|B|^2)\cQ(\Phi,h_{11}^{0})\notag\\
&+2A|A|^2\cQ(\overline\Phi,h_{20}^{+})\notag\\
&+2A|B|^2\cQ(\overline\Phi,h_{20}^{0})\notag\\
&+2A|B|^2\cQ(\Phi,h_{11}^{+}).
\label{eq:Q-feedback-resonant}
\end{align}
It is useful to collect the self- and cross-interaction vectors:
\begin{subequations}
\label{eq:G-vectors}
\begin{align}
G_{\mathrm{s}}={}&
  2\cQ(\Phi,h_{11}^{0})
 +2\cQ(\overline\Phi,h_{20}^{+})
 +3\cC(\Phi,\Phi,\overline\Phi),
\label{eq:Gs}\\
G_{\mathrm{c}}={}&
  2\cQ(\Phi,h_{11}^{0})
 +2\cQ(\overline\Phi,h_{20}^{0})
 +2\cQ(\Phi,h_{11}^{+})
 +6\cC(\Phi,\Phi,\overline\Phi).
\label{eq:Gc}
\end{align}
\end{subequations}
Therefore the full $E_+$ vector coefficient in \eqref{eq:F3-unforced} is
\begin{equation}
  -A_\tau\Phi+\mu A M_1(k_0)\Phi
  +A|A|^2G_{\mathrm{s}}+A|B|^2G_{\mathrm{c}}.
  \label{eq:F3-Eplus}
\end{equation}
Reflection symmetry produces the $E_-$ coefficient
\begin{equation}
  -B_\tau\Phi+\mu B M_1(k_0)\Phi
  +B|B|^2G_{\mathrm{s}}+B|A|^2G_{\mathrm{c}}.
  \label{eq:F3-Eminus}
\end{equation}

\section{Unforced amplitude equations}
\label{sec:normal-form}

Apply $\Psi^*$ to \eqref{eq:F3-Eplus} and
\eqref{eq:F3-Eminus}, use $\Psi^*\Phi=1$, and define
\begin{equation}
  a:=\mu\Psi^*M_1(k_0)\Phi,
  \qquad
  b:=\Psi^*G_{\mathrm{s}},
  \qquad
  c:=\Psi^*G_{\mathrm{c}}.
  \label{eq:abc}
\end{equation}
The Turing--Hopf normal form is
\begin{subequations}
\label{eq:normal-form}
\begin{align}
  A_\tau&=A\left(a+b|A|^2+c|B|^2\right),
  \label{eq:A-normal}\\
  B_\tau&=B\left(a+b|B|^2+c|A|^2\right).
  \label{eq:B-normal}
\end{align}
\end{subequations}
The coefficients are generally complex.  Their imaginary parts correct the
temporal phases/frequencies; their real parts determine amplitudes and
stability.

Writing $A=r_+\ee^{\ii\theta_+}$ and
$B=r_-\ee^{\ii\theta_-}$ yields
\begin{subequations}
\begin{align}
  (r_+)_\tau&=r_+\left(a_r+b_r r_+^2+c_r r_-^2\right),\\
  (r_-)_\tau&=r_-\left(a_r+b_r r_-^2+c_r r_+^2\right),
\end{align}
\end{subequations}
where $a_r=\RePart(a)$, $b_r=\RePart(b)$, and $c_r=\RePart(c)$.

For $a_r>0$ just beyond onset:
\begin{itemize}
  \item A traveling wave has $(r_+,r_-)=(\rho,0)$ or $(0,\rho)$ with
  \begin{equation}
    \rho^2=-\frac{a_r}{b_r}.
  \end{equation}
  It exists supercritically if $b_r<0$ and is stable to the opposite traveling
  mode if $c_r-b_r<0$.

  \item A standing wave has $r_+=r_-=\rho$ with
  \begin{equation}
    \rho^2=-\frac{a_r}{b_r+c_r}.
  \end{equation}
  It exists supercritically if $b_r+c_r<0$ and is stable to amplitude
  imbalance if $c_r-b_r>0$.
\end{itemize}
To leading order, the physical patterns are
\begin{align}
  u_{\mathrm{TW}}&\simeq
  2\epsilon\rho\RePart\!\left(
  \Phi\ee^{\ii(\omega_0t\pm k_0x+\theta)}\right),\\
  u_{\mathrm{SW}}&\simeq
  4\epsilon\rho\cos(k_0x+\vartheta_x)
  \RePart\!\left(\Phi\ee^{\ii(\omega_0t+\vartheta_t)}\right).
\end{align}

\section{External forcing and resonance selection}
\label{sec:forcing}

\subsection{A bookkeeping rule that retains every term}

Associate a fast frequency--wavenumber pair with each first-order factor:
\begin{equation}
\begin{array}{c|cccccc}
\text{factor}&A& B&\overline A&\overline B&P&S\\ \hline
(\Omega,h)&(\omega_0,k_0)&(\omega_0,-k_0)&
(-\omega_0,-k_0)&(-\omega_0,k_0)&(\omega,0)&(0,k)
\end{array}
\label{eq:harmonic-table}
\end{equation}
Conjugate forcing factors have pairs $(-\omega,0)$ and $(0,-k)$.
Multiplication adds the pairs.  A term is resonant in the $A$ equation exactly
when its sum is $(\omega_0,k_0)$, and it is resonant in the $B$ equation
exactly when its sum is $(\omega_0,-k_0)$.

There is also an amplitude-order rule.  Each critical factor costs one power
of $\epsilon$, each $P$ or $\overline P$ costs $n$ powers, and each $S$ or
$\overline S$ costs $m$ powers.  A term competes with the cubic normal form
when the total cost is three.

For computer algebra or a model-specific calculation, write
$\widehat u_j(\Omega,h)$ for the coefficient of
$\ee^{\ii(\Omega s+hx)}$ in $u_j$.  For every noncritical harmonic,
\begin{align}
\widehat u_2(q)=R(q)\bigg[&
  \sum_{q_1+q_2=q}
  \cQ\!\left(\widehat u_1(q_1),\widehat u_1(q_2)\right)
  +\widehat F_2(q)\bigg],
\label{eq:harmonic-u2}
\end{align}
where $q=(\Omega,h)$ and $R(q)=R(\Omega,h)$.  At a critical harmonic
$q_c=(\omega_0,\pm k_0)$, the complete third-order vector is
\begin{align}
\widehat F^{(3)}(q_c)={}&
  -\partial_\tau\widehat u_1(q_c)
  +\mu M_1(\pm k_0)\widehat u_1(q_c)+\widehat F_3(q_c)
  \notag\\
&+\sum_{q_1+q_2=q_c}\left[
  \cQ(\widehat u_1(q_1),\widehat u_2(q_2))
 +\cQ(\widehat u_2(q_1),\widehat u_1(q_2))\right]
\notag\\
&+\sum_{q_1+q_2+q_3=q_c}
  \cC(\widehat u_1(q_1),\widehat u_1(q_2),
  \widehat u_1(q_3)).
\label{eq:full-harmonic-F3}
\end{align}
The Fredholm condition is simply $\Psi^*\widehat F^{(3)}(q_c)=0$.
Equations \eqref{eq:harmonic-u2}--\eqref{eq:full-harmonic-F3} are an exhaustive
recipe: they include the many nonresonant forcing products, forcing-induced
mean fields, and coefficient shifts that are easy to miss by hand.

\subsection{Why neither separated forcing is directly resonant}

The temporal forcing has pair $(\omega,0)$.  Even when $\omega=\omega_0$, its
spatial projection onto $\ee^{\pm\ii k_0x}$ vanishes because $k_0\ne0$.
The spatial forcing has pair $(0,k)$.  Even when $k=k_0$, its temporal
projection onto $\ee^{\ii\omega_0s}$ vanishes because $\omega_0\ne0$.
Therefore neither separated forcing appears as a direct additive term in
\eqref{eq:normal-form}.  Nonlinearity must mix harmonics.

For a noncritical external harmonic define its linear response vectors
\begin{equation}
  p_\omega:=R(\omega,0)P,
  \qquad
  s_k^+:=R(0,k)S,
  \qquad
  s_k^-:=R(0,-k)\overline S.
  \label{eq:forced-responses}
\end{equation}

\subsection{Temporal \texorpdfstring{$2{:}1$}{2:1} parametric resonance}

Let the temporal forcing be $O(\epsilon^2)$, so $n=2$, and take
\begin{equation}
  \omega=2\omega_0+\epsilon^2\delta_t.
\end{equation}
Its second-order response is
$p_2\ee^{2\ii\omega_0s}\ee^{\ii\delta_t\tau}$ with
\begin{equation}
  p_2=R(2\omega_0,0)P.
\end{equation}
The products
\begin{equation}
  (2\omega_0,0)+(-\omega_0,k_0)=(\omega_0,k_0),
\end{equation}
and
\begin{equation}
  (2\omega_0,0)+(-\omega_0,-k_0)=(\omega_0,-k_0)
\end{equation}
couple conjugate counter-propagating modes.  The forced normal form acquires
\begin{subequations}
\begin{align}
  A_\tau&=A(a+b|A|^2+c|B|^2)
  +\eta_t\ee^{\ii\delta_t\tau}\overline B,\\
  B_\tau&=B(a+b|B|^2+c|A|^2)
  +\eta_t\ee^{\ii\delta_t\tau}\overline A,
\end{align}
\end{subequations}
where, for symmetric $\cQ$,
\begin{equation}
  \eta_t=2\Psi^*\cQ(\overline\Phi,p_2).
  \label{eq:eta-time}
\end{equation}
Thus temporal forcing at twice the Hopf frequency parametrically locks pairs
of counter-propagating waves.  The exponent $n=2$ is required because the
forcing response multiplies one $O(\epsilon)$ critical mode and must enter at
$O(\epsilon^3)$.

\subsection{Spatial \texorpdfstring{$2{:}1$}{2:1} parametric resonance}

Let the spatial forcing be $O(\epsilon^2)$, so $m=2$, and choose
$k=2k_0$.  Put
\begin{equation}
  s_2^+=R(0,2k_0)S,
  \qquad
  s_2^-=R(0,-2k_0)\overline S.
\end{equation}
Because
\begin{equation}
  (0,2k_0)+(\omega_0,-k_0)=(\omega_0,k_0),
\end{equation}
and the conjugate relation holds for the other direction, the normal form is
\begin{subequations}
\begin{align}
  A_\tau&=A(a+b|A|^2+c|B|^2)+\eta_x^+B,\\
  B_\tau&=B(a+b|B|^2+c|A|^2)+\eta_x^-A,
\end{align}
\end{subequations}
with
\begin{equation}
  \eta_x^+=2\Psi^*\cQ(\Phi,s_2^+),
  \qquad
  \eta_x^-=2\Psi^*\cQ(\Phi,s_2^-).
  \label{eq:eta-space}
\end{equation}
For a real cosine forcing in a reflection-symmetric system, these coefficients
obey the corresponding conjugacy relation.  This term linearly scatters a
right-traveling mode into a left-traveling mode and tends to pin standing
waves to the imposed spatial pattern.

\subsection{Mixed additive spatio-temporal resonance}

The product of the temporal and spatial responses has pairs
$(\omega,\pm k)$.  It is additive at the critical harmonic when
\begin{equation}
  \omega=\omega_0+\epsilon^2\delta,
  \qquad k=k_0.
  \label{eq:mixed-resonance}
\end{equation}
To enter at the same $O(\epsilon^3)$ order as the cubic normal form, its power
must satisfy
\begin{equation}
  n+m=3.
  \label{eq:nplusm}
\end{equation}
For positive natural numbers this gives $(n,m)=(1,2)$ or $(2,1)$.  Define
\begin{equation}
  p_0=R(\omega_0,0)P,
  \qquad
  s_0^+=R(0,k_0)S,
  \qquad
  s_0^-=R(0,-k_0)\overline S.
\end{equation}
The additive terms are
\begin{equation}
  f_+=2\Psi^*\cQ(p_0,s_0^+),
  \qquad
  f_-=2\Psi^*\cQ(p_0,s_0^-).
  \label{eq:mixed-coefficients}
\end{equation}
The normal form therefore contains
\begin{subequations}
\begin{align}
  A_\tau&=A(a+b|A|^2+c|B|^2)
  +f_+\ee^{\ii\delta\tau}+\text{other allowed forcing terms},\\
  B_\tau&=B(a+b|B|^2+c|A|^2)
  +f_-\ee^{\ii\delta\tau}+\text{other allowed forcing terms}.
\end{align}
\end{subequations}
Here ``other allowed forcing terms'' matters when one forcing is
$O(\epsilon)$: two copies of that response together with one critical mode
also cost three powers of $\epsilon$ and can shift the linear coefficient.
Those shifts are generated without omission by
\eqref{eq:harmonic-u2}--\eqref{eq:full-harmonic-F3}.

\subsection{Summary of the leading resonance balances}

\begin{center}
\begin{tabular}{@{}L{0.22\textwidth}L{0.19\textwidth}L{0.24\textwidth}L{0.24\textwidth}@{}}
\toprule
Mechanism & Exponent balance & Frequency/wavenumber condition & Amplitude-equation effect \\
\midrule
Unforced saturation & $1+1+1=3$ & automatic & $A|A|^2$, $A|B|^2$ and reflected terms \\
Parameter variation & $2+1=3$ & critical mode & $aA$, $aB$ \\
Temporal parametric forcing & $n+1=3$, hence $n=2$ & $\omega=2\omega_0+O(\epsilon^2)$ & $\overline B$ in the $A$ equation and $\overline A$ in the $B$ equation \\
Spatial parametric forcing & $m+1=3$, hence $m=2$ & $k=2k_0$ & $B$ in the $A$ equation and $A$ in the $B$ equation \\
Mixed additive forcing & $n+m=3$ & $\omega=\omega_0+O(\epsilon^2)$, $k=k_0$ & constants $f_\pm\ee^{\ii\delta\tau}$ \\
Forcing-induced linear shift & $2n+1=3$ or $2m+1=3$ & forcing and its conjugate have zero net pair & $|P|^2A$ for $n=1$, or $|S|^2A$ for $m=1$, and reflected terms \\
\bottomrule
\end{tabular}
\end{center}

More complicated resonances require no new principle.  Form every allowed
quadratic or cubic monomial, add its frequency--wavenumber pairs from
\eqref{eq:harmonic-table}, check whether the result is
$(\omega_0,\pm k_0)$, and check whether its $\epsilon$ cost is three.  If a
resonant term appears already at order two, the $\tau=\epsilon^2t$ scaling is
inconsistent; the forcing must be weakened or a faster slow time must be
introduced.

\section{Specialization to the excitatory--inhibitory neural field}

For the neural field in the proposal, write $U=(u,v)^T$ and linearize about a
spatially uniform equilibrium.  Let $f_e^{(r)}$ and $f_i^{(r)}$ denote the
$r$th derivatives of the firing-rate function evaluated at the excitatory and
inhibitory equilibrium inputs.  With Fourier coefficients
$\widehat W_{ab}(h)$, the linear Fourier matrix is
\begin{equation}
M_0(h)=
\begin{pmatrix}
\dfrac{-1+f_e'\widehat W_{ee}(h)}{\tau_e}
&\dfrac{-f_e'\widehat W_{ie}(h)}{\tau_e}\\[6pt]
\dfrac{f_i'\widehat W_{ei}(h)}{\tau_i}
&\dfrac{-1-f_i'\widehat W_{ii}(h)}{\tau_i}
\end{pmatrix}.
\label{eq:neural-M0}
\end{equation}
For fields $U=(u,v)^T$ and $V=(\widetilde u,\widetilde v)^T$, define
\begin{align}
\ell_e(U)&=W_{ee}*u-W_{ie}*v,
&\ell_i(U)&=W_{ei}*u-W_{ii}*v.
\end{align}
Then the polarized Taylor terms are
\begin{equation}
\cQ(U,V)=
\begin{pmatrix}
\dfrac{f_e''}{2\tau_e}\ell_e(U)\ell_e(V)\\[5pt]
\dfrac{f_i''}{2\tau_i}\ell_i(U)\ell_i(V)
\end{pmatrix},
\label{eq:neural-Q}
\end{equation}
and
\begin{equation}
\cC(U,V,W)=
\begin{pmatrix}
\dfrac{f_e'''}{6\tau_e}\ell_e(U)\ell_e(V)\ell_e(W)\\[5pt]
\dfrac{f_i'''}{6\tau_i}\ell_i(U)\ell_i(V)\ell_i(W)
\end{pmatrix}.
\label{eq:neural-C}
\end{equation}
Equations \eqref{eq:neural-M0}--\eqref{eq:neural-C}, together with the choice
of bifurcation parameter that determines $M_1(h)$, provide every object needed
in \eqref{eq:h-vectors}, \eqref{eq:G-vectors}, and \eqref{eq:abc}.

\section{Model-specific computation checklist}

For a concrete parameter set, the calculation can be completed in the
following order.
\begin{enumerate}[label=\arabic*.]
  \item Compute $M_0(h)$ and scan integer $h$ for a simple pair
  $\pm\ii\omega_0$ crossing at $h=\pm k_0$.
  \item Compute $\Phi$ and $\Psi$ from \eqref{eq:eigenvectors}, then normalize
  $\Psi^*\Phi=1$.
  \item Verify the transversality condition \eqref{eq:transversality} and the
  nonresonance/invertibility conditions \eqref{eq:nonresonance-matrices}.
  \item Evaluate $q_{20}$ and $q_{11}$ from \eqref{eq:q20q11}.
  \item Solve the six small linear systems in \eqref{eq:h-vectors}; explicit
  matrix inversion is unnecessary and numerically less stable.
  \item Form $G_{\mathrm{s}}$ and $G_{\mathrm{c}}$ from
  \eqref{eq:G-vectors}, then project with $\Psi^*$ to obtain $a,b,c$.
  \item Compare $\RePart(b)$ and $\RePart(c)$ using the criteria in
  Section \ref{sec:normal-form}.
  \item Add the desired forcing harmonics using
  \eqref{eq:harmonic-u2}--\eqref{eq:full-harmonic-F3}.  The resonance table
  determines which choices of $n,m,\omega,k$ affect the leading amplitude
  equations.
\end{enumerate}

\section{What cannot be made numerical yet}

The derivation above is complete at the abstract level, but the following data
are required for numerical coefficients:
\begin{itemize}
  \item the exact entries of $M_0(h)$ and $M_1(h)$, or equivalently all neural
  kernels, time constants, equilibrium values, and the parameter being varied;
  \item the components and phases of the forcing vectors $P$ and $S$;
  \item the intended temporal and spatial resonance, including any detuning;
  \item the convention for Fourier coefficients of the convolution kernels.
\end{itemize}
Once these are supplied, all remaining work consists of finite-dimensional
linear solves and inner products listed in the checklist.

\begin{thebibliography}{9}

\bibitem{curtu-ermentrout}
R. Curtu and B. Ermentrout,
``Pattern Formation in a Network of Excitatory and Inhibitory Cells with
Adaptation,''
\emph{SIAM Journal on Applied Dynamical Systems}, 3(3), 191--231, 2004.
\href{https://doi.org/10.1137/030600503}{doi:10.1137/030600503}.

\end{thebibliography}

\end{document}
